\subsection{Entanglement as a Resolution Event in Reprojection Dynamics}
  \label{subsec:entanglement-critical-compression}

  The preceding sections establish the structural origin of non-factorizable
  correlations.
  This section gives a unified account of entanglement, disentanglement,
  anti-correlated entanglement, entanglement swapping, and monogamy within a
  single mechanism: the finite-resolution reprojection sequence.

  \paragraph{Generic reprojection.}
    At each step $t \to t+1$ the effective state $\chi_{\mathrm{eff}}(t)$ is
    reprojected through the fixed mapping $\Pi$.
    The generic outcome is a deformed image: information is lost
    (non-injectivity), and the new effective state differs from the old one by
    projective noise.
    Two excitations $A$ and $B$ that $\Pi$ can resolve as distinct produce
    separable fibers $\Pi^{-1}(o_A) \times \Pi^{-1}(o_B)$, and no entanglement
    arises.

  \paragraph{Entanglement.}
    Entanglement arises when the geometry of $\chi_{\mathrm{eff}}(t)$ places
    two degrees of freedom $A$ and $B$ within the same resolution cell of $\Pi$
    in the relevant observable sector.
    The joint fiber $\Pi^{-1}(o_{AB})$ is non-decomposable, and the global
    balance invariant $I_{AB}$ is infra-projectable: it cannot be assigned
    independently to $A$ and $B$.
    No pre-selected or distinguished configuration of $\chi$ is involved.

  \paragraph{Direct and anti-correlated entanglement.}
    The sign of the correlation is not a separate mechanism but a property of
    the infra-projectable invariant $I_{AB}$ carried by the fiber.
    When the invariant selects decompositions in which local values coincide
    (e.g.\ $s_A = s_B = +\tfrac{1}{2}$), the correlation is direct.
    When the invariant is a conservation law that forces opposite local values
    (e.g.\ $s_A + s_B = 0$), the correlation is anti-correlated.
    The singlet and triplet states correspond to distinct global invariants
    carried by distinct fibers, not to distinct mechanisms.

  \paragraph{Disentanglement by measurement or decoherence.}
    Measurement and environmental coupling modify the effective state
    $\chi_{\mathrm{eff}}(t)$.
    The subsequent reprojection places $A$ and $B$ at resolvable positions in
    the observable sector: the joint fiber becomes decomposable, the global
    invariant $I_{AB}$ becomes locally projectable, and the entanglement
    disappears.
    No collapse or destruction of information is involved.
    The balance constraint that was infra-projectable becomes projectable once
    the effective state has moved sufficiently in the space of observable
    configurations; $\Pi$ has not changed, only the image it is reading.

  \paragraph{Entanglement swapping.}
    Suppose at step $t_1$ excitations $A$ and $B$ fall below resolution and
    share the infra-projectable invariant $I_{AB}$; at step $t_2$ excitations
    $B$ and $C$ share $I_{BC}$.
    A measurement on $B$ forces a locally admissible decomposition of $I_{AB}$
    and simultaneously restricts the admissible decompositions of $I_{BC}$.
    The residual constraint propagates to $A$ and $C$: the subsequent effective
    state places $A$ and $C$ in a jointly constrained fiber $\Pi^{-1}(o_{AC})$,
    even though $A$ and $C$ have never shared a resolution cell directly.
    Swapping is transitive propagation of fiber constraints through an
    intermediate reprojection event, not a transmission of influence.

  \paragraph{Monogamy.}
    A resolution cell has bounded capacity, set by the Born--Infeld saturation
    condition on the projection fiber.
    A single cell cannot simultaneously carry independent infra-projectable
    invariants for multiple disjoint pairs without saturating.
    Monogamy of entanglement is therefore a capacity constraint on the fiber,
    not an additional postulate.

  \paragraph{Summary.}
    The five phenomena---entanglement, direct and anti-correlated correlations,
    disentanglement, swapping, and monogamy---share a single mechanism: the
    interaction of the evolving effective state $\chi_{\mathrm{eff}}(t)$ with
    the fixed finite-resolution projection $\Pi$.
    Entanglement is a local resolution event in the reprojection dynamics; it
    requires no special configuration of $\chi$, no modification of $\Pi$, and
    no non-local dynamical influence.
    The balance constraints that become infra-projectable originate in the
    admissibility structure of $\chi$---which is a coherence space, not a data
    store---and not in any classically rich hidden state or pre-selected
    substrate configuration.
    This position is distinct from both na\"{\i}ve realism (no observable values
    are stored in $\chi$) and eliminativism (the admissibility structure of
    $\chi$ is the genuine source of all invariants).
    The key principle is: \emph{what lies below the projective resolution is
not a value, but a constraint.}
